\documentclass[12pt,a4paper]{article}
\usepackage[utf8]{inputenc}
\usepackage[T1]{fontenc}
\usepackage[ngerman]{babel}
\usepackage{geometry}
\usepackage{amsmath}
\usepackage{graphicx}
\usepackage{hyperref}
\usepackage{setspace}
\usepackage{csquotes}   % Für bessere Anführungszeichen
\geometry{margin=2.5cm}
\onehalfspacing
\hypersetup{
	colorlinks=true,
	linkcolor=blue,
	urlcolor=blue,
	citecolor=black
}
\begin{document}
	\title{Wissenschaftliche Perspektiven auf Unterschiede zwischen Frauen und Männern\\
		\large{Eine interdisziplinäre Quellenanalyse}}
	\author{}
	\date{\today}
	\maketitle
	\begin{abstract}
		Die Frage nach Unterschieden zwischen den Geschlechtern ist Gegenstand intensiver wissenschaftlicher Debatten. Diese Arbeit bietet eine systematische Übersicht über aktuelle Forschungsergebnisse zu biologischen, psychologischen und sozialen Geschlechterdifferenzen. Basierend auf einer strukturierten Quellenrecherche werden Erkenntnisse aus Genetik, Neurobiologie, Hormonforschung, Kognitionspsychologie, Medizin und Soziologie dargestellt. Die Ergebnisse zeigen, dass Geschlechtsunterschiede in fast allen Domänen gering ausfallen, stark von Umweltfaktoren abhängen und häufig Wechselwirkungen zwischen biologischen und sozialen Einflüssen widerspiegeln. Die Arbeit unterstreicht die Bedeutung der Unterscheidung zwischen biologischem Sex und sozialem Gender sowie die Notwendigkeit geschlechtsspezifischer Ansätze in Forschung und Klinik.
	\end{abstract}
	
	\section{Einführung}
	Die wissenschaftliche Auseinandersetzung mit Geschlechterunterschieden ist von komplexen Debatten geprägt, die oft zwischen biologischen und sozialen Erklärungsansätzen oszillieren. Diese Spannung spiegelt sich in der Forschungsliteratur wider, in der die Ergebnisse je nach Fachdisziplin und methodischem Ansatz variieren. Ziel dieses Artikels ist es, eine ausgewogene, evidenzbasierte Übersicht über die aktuellen wissenschaftlichen Erkenntnisse zu Unterschieden zwischen Frauen und Männern zu geben. Dabei werden folgende zentrale Fragen adressiert: In welchen Bereichen zeigen sich gesicherte Geschlechterdifferenzen? Welche biologischen, psychologischen und sozialen Mechanismen liegen diesen zugrunde? Und wie lassen sich die oft widersprüchlichen Befunde unterschiedlicher Forschungsrichtungen integrieren? Um diese Fragen zu beantworten, wurde eine strukturierte Quellenrecherche in führenden wissenschaftlichen Datenbanken durchgeführt, wobei der Schwerpunkt auf Metaanalysen, systematischen Reviews und aktuellen Primärstudien aus den Jahren 2020 bis 2025 liegt.
	
	\section{Biologische Grundlagen}
	\subsection{Genetische und hormonelle Unterschiede}
	Die biologischen Unterschiede zwischen den Geschlechtern haben ihre Grundlage in den Geschlechtschromosomen. Frauen besitzen typischerweise zwei X-Chromosomen, während Männer ein X- und ein Y-Chromosom aufweisen, wobei das Y-Chromosom das geschlechtsbestimmende SRY-Gen trägt. Diese genetische Disposition beeinflusst nicht nur die Entwicklung primärer und sekundärer Geschlechtsmerkmale, sondern auch die Expression zahlreicher Gene im gesamten Organismus. So verfügen Frauen beispielsweise über zwei X-Chromosomen mit Tumorsuppressorgenen, während Männer bei Funktionsverlust durch Mutation keine zweite Genkopie zur Kompensation haben, was zu Unterschieden in der Krebsanfälligkeit beiträgt. \cite{Hasenburg2025}\\
	Hormonelle Unterschiede sind ebenfalls zentral: Testosteron, das bei Männern in deutlich höheren Konzentrationen vorkommt, beeinflusst Muskelwachstum, Körperbehaarung und Stimmbruch sowie Verhaltensparameter wie Aggressivität und räumliches Denken. Östrogene hingegen, die bei Frauen während des Menstruationszyklus stark fluktuieren, regulieren unter anderem den weiblichen Reproduktionstrakt und das Fettgewebe. Bemerkenswert ist jedoch, dass die Effekte dieser Hormone nicht deterministisch sind. Die Forschung zeigt, dass Hormone in komplexer Weise mit Umweltfaktoren interagieren und dass sogar kurzfristige Veränderungen des Hormonspiegels -- wie sie etwa im weiblichen Zyklus auftreten -- strukturelle und funktionelle Unterschiede im Gehirn modulieren können. \cite{Valk2024}
	
	\subsection{Neuroanatomische und hirnfunktionelle Differenzen}
	Ein besonders intensiv erforschtes Feld betrifft die Frage nach Unterschieden im Gehirn von Frauen und Männern. Metaanalytische Befunde zeigen, dass das durchschnittliche Gehirnvolumen bei Männern etwa 10\% größer ist, was jedoch weitgehend auf den Unterschied in der Körpergröße zurückzuführen ist. Die Relation von Hirnmasse zu Körpermasse ist bei Frauen sogar tendenziell etwas günstiger. \cite{Spektrum} Wichtiger als absolute Größenunterschiede sind jedoch strukturelle und funktionelle Variationen.
	Aktuelle neurowissenschaftliche Studien deuten darauf hin, dass die Gehirne von Frauen und Männern unterschiedliche Netzwerktopologien aufweisen: Männliche Gehirne zeigen demnach eine komplexere hierarchische Organisation, während weibliche Gehirne eine stärkere Modularität und eine höhere Konnektivität innerhalb von Modulen aufweisen. \cite{PMC11507829} Diese Unterschiede sind vor allem in sensorischen und Assoziationsregionen lokalisiert und bleiben auch nach Kontrolle der Gehirngröße bestehen. \cite{NatureCommunications2024} Entscheidend ist jedoch, dass diese Differenzen nicht als \glqq besser\grqq\ oder \glqq schlechter\grqq\ interpretiert werden dürfen, sondern unterschiedliche, aber gleichwertige Verarbeitungsstrategien repräsentieren.
	Eine bemerkenswerte Studie aus dem Jahr 2024 konnte zeigen, dass sich biologische und soziale Geschlechtsaspekte bereits im kindlichen Gehirn in unterschiedlichen Mustern widerspiegeln: Während somatomotorische, visuelle, Kontroll- und limbische Netzwerke vor allem mit dem biologischen Geschlecht assoziiert sind, sind die neuronalen Korrelate von Gender über den gesamten Kortex verteilt. \cite{Science2024} Dies unterstreicht die Notwendigkeit, biologische und soziale Einflüsse getrennt zu betrachten.
	
	\section{Psychologische und kognitive Unterschiede}
	\subsection{Kognitive Fähigkeiten und Verhaltenstendenzen}
	Die populäre Vorstellung, dass Männer und Frauen in ihren kognitiven Fähigkeiten grundlegend verschieden seien, wird durch die empirische Forschung nur sehr eingeschränkt bestätigt. Eine großangelegte Metaanalyse von Janet Hyde, die über 7.000 Einzelstudien umfasste, kommt zu dem Schluss, dass die überwältigende Mehrheit psychologischer Geschlechtsunterschiede gering bis vernachlässigbar ist. \cite{DWDS} Signifikante Unterschiede finden sich lediglich in wenigen Bereichen: So zeigen Männer im Durchschnitt eine bessere mentale Rotation (räumliches Vorstellungsvermögen), während Frauen Vorteile in verbalen Flüssigkeitsaufgaben und bei der Erkennung von Gesichtern und Emotionen haben. Auch bei Aggressivität, Masturbationshäufigkeit und der Ablehnung von Gelegenheitssex zeigen sich konsistente Differenzen. \cite{DWDS}
	Eine aktuelle systematische Übersichtsarbeit und Metaanalyse zum Gender-Equality-Paradoxon liefert hierzu hochrelevante Erkenntnisse: Die Autoren fanden, dass Geschlechtsunterschiede in Persönlichkeit, verbalen Fähigkeiten, episodischem Gedächtnis und negativen Emotionen in Ländern mit höherem Wohlstandsniveau ausgeprägter sind. Demgegenüber fallen Unterschiede im Sexualverhalten, in Partnerpräferenzen und in Mathematik in wohlhabenderen Ländern geringer aus. \cite{Herlitz2025} Dieses auf den ersten Blick paradoxe Muster deutet darauf hin, dass verbesserte Lebensbedingungen und größere Wahlfreiheit dazu führen können, dass angeborene Präferenzen und Neigungen stärker zum Tragen kommen. Es widerlegt die einfache Vorstellung, dass soziale Gleichstellung automatisch zu einer Angleichung psychologischer Profile führt.
	
	\subsection{Intelligenz und kognitive Fähigkeiten}
	Die Frage nach signifikanten Geschlechtsunterschieden in der allgemeinen Intelligenz ist Gegenstand intensiver Forschung. Die überwiegende Evidenz zeigt, dass es keinen signifikanten Unterschied in der globalen Intelligenz (g-Faktor) zwischen Männern und Frauen gibt – sie erreichen im Durchschnitt praktisch identische Gesamt-IQ-Werte. Eine groß angelegte Metaanalyse von Giofrè et al. (2022) mit über 46.000 Teilnehmern zeigte, dass Unterschiede in der fluiden Intelligenz vernachlässigbar sind. Interessanterweise sind Geschlechterunterschiede bei neueren Testversionen (z. B. WISC-V) geringer ausgeprägt, was auf abnehmende Stereotypeffekte hindeuten könnte. \cite{Giofre2022}
	
	Dennoch zeigen sich konsistente, wenn auch geringe Unterschiede in spezifischen kognitiven Profilen: Männer schneiden im Durchschnitt besser bei visuell-räumlichen Aufgaben und bei quantitativem Wissen ab, während Frauen bei der Verarbeitungsgeschwindigkeit und beim Kurzzeitgedächtnis Vorteile haben. Eine weitere wichtige Erkenntnis betrifft die intraindividuelle Variabilität: In manchen Domänen zeigen Männer eine höhere Streuung der Leistungen, was bedeutet, dass es mehr Männer sowohl im unteren als auch im oberen Leistungsbereich gibt. \cite{Giofre2024} Eine gesonderte Analyse an Kindern mit Lernstörungen bestätigt, dass auch in klinischen Stichproben keine globalen Intelligenzunterschiede zwischen den Geschlechtern bestehen. \cite{Toffalini2025}
	
	Der Flynn-Effekt – der langfristige Anstieg der IQ-Werte in der Bevölkerung – verläuft bei Männern und Frauen weitgehend parallel. Zwar schneiden Männer in einzelnen Subskalen höher ab, die generationsübergreifenden Zuwächse sind jedoch geschlechtsunabhängig. \cite{Gnambs2011}
	
	\subsection{Emotionale Verarbeitung und psychische Gesundheit}
	Im Bereich der Emotionsregulation und psychischen Gesundheit zeigen sich ebenfalls interessante Geschlechterunterschiede. Frauen weisen eine höhere Rate internalisierender Störungen wie Depressionen und Angststörungen auf, während Männer häufiger externalisierende Verhaltensweisen und Substanzabhängigkeit zeigen. Eine systematische Übersichtsarbeit aus dem Jahr 2025 untersuchte den Zusammenhang zwischen unerwünschten Kindheitserfahrungen und psychischen Erkrankungen und fand, dass die Effekte je nach Art und kumulativer Exposition der traumatischen Ereignisse variieren. \cite{PsychMeta2025} Frauen sind nicht nur häufiger von Depressionen betroffen, sondern zeigen auch eine höhere Prävalenz von Autoimmunerkrankungen, während Männer ein höheres Risiko für Parkinson und Suizid aufweisen. \cite{KISpotlight}
	Interessant ist auch der Befund einer Metaanalyse zur interozeptiven Genauigkeit (der Fähigkeit, körpereigene Signale wahrzunehmen): Überraschenderweise haben Frauen tendenziell eine geringere interozeptive Genauigkeit als Männer, schätzen ihre emotionale Kompetenz jedoch höher ein und zeigen eine bessere Emotionsverarbeitung. \cite{ScienceDirect2025} Dies deutet auf eine Dissoziation zwischen objektiver Wahrnehmungsfähigkeit und subjektivem Emotionserleben hin, die geschlechtsspezifisch ausgeprägt sein könnte.
	
	\section{Gesundheitliche Unterschiede}
	\subsection{Morbidität, Mortalität und Krankheitsrisiken}
	Männer und Frauen unterscheiden sich auch in signifikanter Weise hinsichtlich ihrer Gesundheitsparameter. Ein zentrales Phänomen ist die sogenannte \glqq morbidity-mortality paradox\grqq: Frauen haben eine höhere Lebenserwartung, verbringen jedoch mehr Lebensjahre in schlechter Gesundheit. \cite{OECD2025} So sind Frauen um 25\% mehr Lebensjahre in schlechter Gesundheit als Männer. \cite{BerlinScienceSurvey}\\
	Die zugrundeliegenden Ursachen sind multifaktoriell und umfassen sowohl biologische als auch Verhaltensfaktoren. So haben Frauen ein höheres Risiko für Autoimmunerkrankungen, Schilddrüsenerkrankungen und Depressionen, während Männer häufiger an Krebs erkranken und eine höhere Mortalität aufweisen. \cite{NatureDiseases2024} Ein wichtiger Faktor ist hierbei das unterschiedliche Risikoverhalten: Männer konsumieren häufiger Alkohol und Nikotin und nehmen seltener Vorsorgeuntersuchungen wahr. \cite{Hasenburg2025}\\
	Besonders augenfällig sind die Unterschiede im Herz-Kreislauf-Bereich: Während Männer insgesamt häufiger an Herzerkrankungen leiden, haben Frauen ein höheres Risiko, an einem Herzinfarkt zu sterben. Zudem sind Frauen biologisch vulnerabler für Nikotin -- Rauchen erhöht das Herzinfarktrisiko bei Frauen um das Siebenfache, bei Männern um das Drei- bis Vierfache. \cite{IASGruppe}
	
	\subsection{Gendermedizin: Konsequenzen für Diagnostik und Therapie}
	Die geschlechtsspezifischen Unterschiede in der Pathophysiologie, Symptompräsentation und Medikamentenwirkung haben weitreichende klinische Implikationen. Ein zentrales Problem ist der historische Bias in der medizinischen Forschung: Therapieleitlinien basieren überwiegend auf männlichen Studienpopulationen, da Frauen in klinischen Studien lange Zeit unterrepräsentiert waren. \cite{NatureDiseases2024} Dies führt zu einer erheblichen Wissenslücke hinsichtlich geschlechtsspezifischer Wirkmechanismen und optimaler Dosierungen.\\
	Konkret zeigt sich, dass Frauen unter zielgerichteten Immun- und Chemotherapien ein höheres Risiko für schwerere Nebenwirkungen haben. \cite{Hasenburg2025} Die Dosierung erfolgt bislang nach Körperoberfläche, nicht jedoch nach den unterschiedlichen Körperzusammensetzungen oder der unterschiedlichen Clearance von Medikamenten. Zudem werden Krankheitssymptome bei Frauen häufiger übersehen oder fehlgedeutet -- so berichten betroffene Frauen bei Herzinfarkten neben Brustschmerzen auch unspezifische Symptome wie Atemnot oder Erschöpfung, die nicht immer als Herzinfarkt erkannt werden. \cite{ClinicBarcelona2025}\\
	Die Gendermedizin hat es sich daher zur Aufgabe gemacht, diese Lücken zu schließen. Sie untersucht systematisch die Unterschiede zwischen Männern und Frauen in Diagnostik, Prävention und Therapie und entwickelt geschlechtsspezifische Ansätze. \cite{BMWFW2025} Erste Erfolge sind sichtbar, dennoch bleibt der Weg zu einer umfassend individualisierten, geschlechtssensiblen Medizin noch weit.
	
	\section{Soziale und kulturelle Einflüsse}
	\subsection{Geschlechterrollen und Sozialisation}
	Biologische Unterschiede zwischen den Geschlechtern sind stets in einen sozialen und kulturellen Kontext eingebettet. Geschlechterrollen -- definiert als gesellschaftlich konstruierte Normen, Erwartungen und Werte, die mit einem bestimmten Geschlecht assoziiert werden -- beeinflussen maßgeblich das Denken, Fühlen und Handeln von Individuen. \cite{HLS2024} Sie sind nicht biologisch determiniert, sondern werden durch Sozialisationsprozesse von Kindheit an vermittelt und verfestigt. \cite{FES2024}\\
	Die Forschung zeigt, dass bereits im Alter von zwei bis drei Jahren Mädchen in ihren sprachlichen und motorischen Fähigkeiten weiter sind als Jungen. \cite{Zeit2024} Erzieherinnen und Erzieher nehmen Mädchen zudem als kreativer, geduldiger und sozialer wahr, während Jungen später eingeschult werden. Solche frühen Unterschiede können durch die unterschiedliche Behandlung und Erwartungshaltung verstärkt werden und sich verfestigen. Entscheidend ist jedoch, dass Geschlechterrollen einem historischen Wandel unterliegen und sich mit den gesellschaftlichen Bedingungen verändern. \cite{DIW2025}
	
	\subsection{Interaktion von Biologie und Umwelt}
	Die moderne Geschlechterforschung hat das simple Dichotomie-Modell von \glqq Natur versus Erziehung\grqq\ längst hinter sich gelassen. Stattdessen wird heute von einer komplexen Interaktion zwischen biologischen Faktoren (Genetik, Hormone, Neurobiologie) und sozialen Umwelteinflüssen ausgegangen. \cite{Asendorpf2015} Dies zeigt sich besonders deutlich in der bereits erwähnten Gender-Equality-Studie, die aufzeigt, dass Unterschiede in Persönlichkeitsmerkmalen und kognitiven Fähigkeiten in Ländern mit höherem Wohlstand tendenziell größer sind. \cite{Herlitz2025}\\
	Die Interaktion kann auch auf neurobiologischer Ebene nachgewiesen werden: Eine aktuelle Studie im Fachjournal \glqq Science Advances\grqq\ konnte zeigen, dass sowohl das biologische Geschlecht (sex) als auch das soziale Geschlecht (gender) mit jeweils distinkten Mustern funktioneller Hirnkonnektivität assoziiert sind. \cite{Science2024} Mit anderen Worten: Die Art und Weise, wie wir unsere Geschlechtsidentität erleben und ausdrücken, hinterlässt nachweisbare Spuren in unserem Gehirn. Diese Befunde belegen, dass eine isolierte Betrachtung biologischer Determinanten nicht ausreicht, um die Vielfalt menschlicher Geschlechterunterschiede zu erklären.
	
	\section{Fazit und Ausblick}
	Die wissenschaftliche Forschung zu Geschlechterunterschieden hat in den letzten Jahrzehnten ein differenziertes Bild gezeichnet, das weit entfernt ist von simplen Stereotypen oder biologistischen Reduktionismen. Zusammenfassend lassen sich folgende Kernbotschaften festhalten:
	\begin{enumerate}
		\item \textbf{Geringe Effektstärken:} In den meisten psychologischen, kognitiven und verhaltensbezogenen Domänen sind die Unterschiede zwischen den Geschlechtern gering bis vernachlässigbar. \cite{DWDS}
		\item \textbf{Keine globalen Intelligenzunterschiede:} Männer und Frauen unterscheiden sich nicht in der allgemeinen Intelligenz (g-Faktor). Lediglich in spezifischen Teilbereichen zeigen sich geringfügige, domänenspezifische Differenzen. \cite{Giofre2022}
		\item \textbf{Biologische Grundlagen:} Es gibt zweifellos biologische Unterschiede -- genetisch, hormonell und neuroanatomisch. Ihre Auswirkungen sind jedoch weitaus kleiner und vor allem variabler, als lange Zeit angenommen.
		\item \textbf{Kontextabhängigkeit:} Die Stärke und Richtung von Geschlechterdifferenzen variiert mit den kulturellen, sozialen und ökonomischen Rahmenbedingungen. \cite{Herlitz2025}
		\item \textbf{Klinische Relevanz:} In der Medizin sind geschlechtsspezifische Unterschiede von großer Bedeutung. Die Gendermedizin ist ein wachsendes Feld, das beachtliche Verbesserungen in Diagnostik und Therapie verspricht. \cite{NatureDiseases2024}
		\item \textbf{Interaktionsperspektive:} Anstatt nach einer monokausalen Erklärung zu suchen, ist eine interaktionsistische Perspektive notwendig, die das komplexe Zusammenspiel von genetischen, hormonellen, neuronalen, psychologischen und soziokulturellen Faktoren berücksichtigt.
	\end{enumerate}
	Diese Erkenntnisse haben weitreichende Implikationen -- sowohl für die wissenschaftliche Praxis als auch für gesellschaftliche Debatten. In der Forschung ist eine konsequente Trennung und gemeinsame Erhebung von Sex und Gender notwendig. Klinische Studien müssen in ausreichendem Maße Frauen einschließen und geschlechtsspezifische Subgruppenanalysen durchführen. Und in der öffentlichen Diskussion gilt es, sowohl die Realität biologischer Unterschiede anzuerkennen als auch die bedeutsame Rolle sozialer Konstruktionsprozesse zu würdigen. Die wissenschaftliche Erforschung der Geschlechterunterschiede ist damit nicht nur ein faszinierendes akademisches Unterfangen, sondern auch eine gesellschaftlich hochrelevante Aufgabe.
	
	\section{Quellenverzeichnis}
	\bibliographystyle{plain}
	\begin{thebibliography}{99}
		\bibitem{Hasenburg2025} Hasenburg, A. (2025). Genderspezifische Therapien und Prävention. \textit{gynäkologie + geburtshilfe}, 30, 56.
		\bibitem{Valk2024} Valk, S. et al. (2024). Relating sex-bias in human cortical and hippocampal microstructure to sex hormones. \textit{Nature Communications}, 15, 7279.
		\bibitem{Spektrum} Hilbig, H. (o.D.). Geschlechtsunterschiede aus neurowissenschaftlicher Sicht. \textit{Lexikon der Neurowissenschaft}. Spektrum.
		\bibitem{PMC11507829} (2024). Sex Differences in Hierarchical and Modular Organization of Functional Brain Networks. \textit{Entropy}, 26(10), 864.
		\bibitem{NatureCommunications2024} (2024). Sex differences in functional cortical organization reflect differences in network topology rather than cortical morphometry. \textit{Nature Communications}, 15, 7714.
		\bibitem{Science2024} (2024). Functional brain networks are associated with both sex and gender in children. \textit{Science Advances}, 10(28), eadn4202.
		\bibitem{DWDS} DWDS (2021). Geschlechterunterschied. In: \textit{Digitales Wörterbuch der deutschen Sprache}.
		\bibitem{Herlitz2025} Herlitz, A., Hönig, I., Hedebrant, K., \& Asperholm, M. (2025). A Systematic Review and New Analyses of the Gender-Equality Paradox. \textit{Perspectives on Psychological Science}, 20(3), 503-539.
		\bibitem{PsychMeta2025} (2025). Gender differences in the associations of adverse childhood experiences with depression and anxiety. \textit{Journal of Affective Disorders}.
		\bibitem{KISpotlight} Karolinska Institutet (2020). Spotlight on sex and gender differences in health.
		\bibitem{ScienceDirect2025} (2025). A systematic review and meta-analysis of sex differences in the relationship between cardiac interoceptive accuracy and emotion. \textit{Neuroscience \& Biobehavioral Reviews}.
		\bibitem{OECD2025} OECD (2025). Gesundheit auf einen Blick 2025 (Auszugsweise ...).
		\bibitem{BerlinScienceSurvey} Berlin Science Survey (o.D.). Frauen sind deutlich stärker belastet als Männer.
		\bibitem{NatureDiseases2024} (2024). Sex difference in human diseases: mechanistic insights and clinical implications. \textit{Signal Transduction and Targeted Therapy}, 9, 238.
		\bibitem{IASGruppe} IAS-Gruppe (o.D.). Die größten Risikofaktoren für die Entwicklung von Herz-Kreislauf-Erkrankungen.
		\bibitem{ClinicBarcelona2025} Clinic Barcelona (2025). Health with a sex and gender perspective.
		\bibitem{BMWFW2025} BMWFW (2025). Nachlese zum Science Talk \glqq Wie Gendern auch Leben ...\grqq.
		\bibitem{HLS2024} HLS (2024). Geschlechterrollen. In: \textit{Historisches Lexikon der Schweiz}.
		\bibitem{FES2024} FES (2024). Geschlechterrollen. \textit{Gender Glossar}. Friedrich-Ebert-Stiftung.
		\bibitem{Zeit2024} Die Zeit (2024). Ausgabe 39 / 2024: Unsere Quellen.
		\bibitem{DIW2025} DIW Berlin (2025). Einstellungen zu Geschlechterrollen sind im Wandel.
		\bibitem{Asendorpf2015} Asendorpf, J. B. (2015). Geschlechtsunterschiede. In: \textit{Persönlichkeitspsychologie für Bachelor}. Springer.
		\bibitem{Giofre2022} Giofrè, D., Allen, K., Toffalini, E., \& Caviola, S. (2022). The Impasse on Gender Differences in Intelligence: A Meta-Analysis on WISC Batteries. \textit{Educational Psychology Review}, 34, 2503–2528.
		\bibitem{Toffalini2025} Toffalini, E., Giofrè, D., \& Cornoldi, C. (2025). Sex Differences in Intelligence on the WISC: A Meta-Analysis on Children with Learning Difficulties. \textit{Journal of Intelligence}, 13(2), 18.
		\bibitem{Giofre2024} Giofrè, D., Toffalini, E., \& Calcagnì, A. (2024). Sex differences in cognition: A meta-analysis of variance ratios. \textit{Personality and Individual Differences}, 226, 112684.
		\bibitem{Gnambs2011} Gnambs, T., \& Stieger, S. (2011). Female Flynn effects: No evidence for sex differences in the secular increase of intelligence. \textit{Intelligence}, 39(5), 319-325.
	\end{thebibliography}
\end{document}